\section{Лекция 11. Теорема о единственности. Теорема Вейерштрасса}
\textbfit{Опр.} Пусть $f(\infty) = 0, \ f\in \mathcal{O}(|z|>R)$. Тогда порядком нуля в бесконечности называют порядок нуля $f\left(\frac{1}{z}\right)$ в 0.

$g(z) = f\left(\frac{1}{z}\right) = z^N \tilde{g}(z), \ \tilde{g}(z) \neq 0$

$f(z) = g \left(\frac{1}{z}\right) = \frac{\tilde{g}\left(\frac{1}{z}\right)}{z^N} = \frac{1}{z^N}g_1(z), \ \tilde{g} \left(\frac{1}{z}\right) \neq 0$

\subsection{Теорема о единственности}
\textbfit{Теорема} (о единственности). Пусть $f, g \in \mathcal{O}(D), f \overset{E}{=} g, \ E \subset D, E $ имеет предельную точку в $D.$ Тогда $f \overset{D}{=} g.$
\[\triangleright \text{ Рассмотрим } F = f - g, \ F \in \mathcal{O}(D), F \overset{E}{=} 0\]
\[\text{Пусть $a$ -- предельная точка } E, a \in D \Rightarrow F(a) = 0 \text{. Тогда $a$ -- неизолированный ноль.} \Rightarrow\]
\[\Rightarrow \exists U_R(a) \subset D: F = 0 \text{ на } U_R(a)\]
\[\text{Пусть } \exists b \in D: F(b) \neq 0 \text{. Рассмотрим путь } \Gamma:[0, 1] \to D: \Gamma(0) = a, \Gamma(1) = b\]
\[\text{Пусть } T = \inf \{t \in [0, 1]: F(\Gamma(t)) \neq 0\} \text{ -- непусто, ограничено}\]
\[T > 0, \text{ т.к. } \exists \delta: \forall t < \delta \Gamma(t) \in U_R(a) \Rightarrow F(\Gamma(t))=0\]
\[T \neq 1 \text{ в силу непрерывности } F \text{ (иначе $F(b) = 0$)}\]
\[0 < T < 1. \text{ По непрерывности } F(\Gamma(T)) = 0 \Rightarrow \Gamma(T) \text{ -- неизолированный ноль} \Rightarrow\]
\[\Rightarrow \exists U_r(\Gamma(T)): \forall z \in U_r(\Gamma(T)) \ F(z) = 0 \Rightarrow \exists \delta_1: \forall t \in (T, T + \delta_1) \ \Gamma(t) \in U_r(\Gamma(T)), F(\Gamma(t)) = 0 \Rightarrow\]
\[\Rightarrow \nexists \inf \text{ -- противоречие }\blacktriangleleft\]

\textbfit{Пр.} $f \in \mathcal{O}(|z|<1), z_n \to \frac{1}{n} \in U_1(0), f\left(\frac{1}{n}\right) = 0 \Rightarrow f \equiv 0$
\[f \in \mathcal{O}(|z-1|<1), z_n \to \frac{1}{n} \in U_1(0), f \left(\frac{1}{n}\right) = 0 \overset{?}{\Rightarrow} f = \sin \frac{\pi}{z}\]
\[f \in \mathcal{O}(\mathbb{C}), z_n \to n \in U_1(0), f(n) = 0 \overset{?}{\Rightarrow} f = \sin \pi z\]

\textbfit{Замечание} Теоремы Лагранжа не существует в комплексном анализе.

\subsection{Теорема Вейерштрасса}
\textbfit{Теорема} (Вейерштрасса). Пусть $f_n \in \mathcal{O}(D)$. Пусть $\forall K \subset D, K$ -- компакт $f_n \overset{K}{\rightrightarrows} f$ (меньше, чем сходимость на всей $D$). Тогда

1) $f \in \mathcal{O}(D)$

2) $\forall z \in D \ \forall k f_n^{(k)}(z) \to f^{(k)}(z)$

3) $\forall K \subset D, K -$ компакт $f_n^{(k)}(z) \overset{K}{\rightrightarrows} f(z)$
\[\triangleright \ 1) \text{ Покажем } f \in \mathcal{O}(D)\]
\[f \in C(D) \text{ (т.к. $C(D) \ni f_n \overset{K}{\rightrightarrows} f \Rightarrow f \in C(K) \forall z_0 \in D \exists \overline{U_R(z_0)} \subset{D} \Rightarrow f \in C(\overline{U_R(z_0)}) \Rightarrow$}\]
\[\Rightarrow f \in C(z_0) \Rightarrow f \in C(D))\]
\[\text{Пусть } \overline{\triangle} \subset D \int_{\partial \overline{\triangle}} f(z) dz = \int_{\partial \overline{\triangle}} (\lim_{n\to\infty} f_n(z))dz \overset{f_n \overset{\partial \triangle}{\rightrightarrows} f}{=} \lim_{n\to\infty} f_n(z) dz = 0, \ f_n(z) = 0\text{ (л. Гурса)} \Rightarrow\]
\[f \in \mathcal{O}(D) \text{ по т. Мореры}\]
\[2) \text{ Пусть $a \in \triangle$. Покажем } \forall k \lim_{n\to\infty} f_n^{(k)}(a) = f^{(k)}(a)\]
\[\exists R > 0: \overline{U_R(a)} \subset D\]
\[\text{Тогда }f^{(k)}(a) \overset{\text{ИФК}}{=} \frac{k!}{2\pi i} \oint_{|z-a|=R} \frac{f(z)}{(z-a)^{k+1}}dz = \frac{k!}{2\pi i} \oint_{|z-a|=R} \left(\lim_{n\to\infty}\frac{f(z)}{(z-a)^{k+1}}\right)dz = (*)\]
\[\sup_{|z-a|=R} \left|\frac{f_n(z)}{(z-a)^{k+1}} - \frac{f(z)}{(z-a)^{k+1}}\right| = \frac{1}{R^{k+1}}\sup_{|z-a|=R} |f_n(z)-f(z)|<\epsilon \forall n > N \text{ т.к. } f_n \overset{|z-a|=R}{\rightrightarrows} f\]
\[(*) \overset{\text{в силу равномерной сходимости}}{=} \lim_{n\to\infty} \frac{k!}{2\pi i} \oint_{|z-a|=R} \frac{f_n(z)}{(z-a)^{k+1}}dz \overset{\text{ИФК}}{=} \lim_{n\to\infty} f_n^{(k)}(a)\]
\[3) \text{ Сначала покажем, что } \forall a \in D \exists R(a): f_n^{(k)} \overset{\overline{U_{R(a)}(a)}}{\rightrightarrows} f^{(k)}(z)\]
\[\exists r(a) : \overline{U_{r(a)}(a)} \subset D. \text{ Возьмём } 0 < R(a) < r(a).\] 
\[\text{ Тогда } \forall z \in \overline{U_{R(a)}} \ f_n^{(k)}(z) - f^{(k)}(z) = \frac{k!}{2\pi i} \oint_{|\zeta-a|=r(a)} \frac{f_n(\zeta)-f(\zeta)}{(\zeta-z)^{k+1}}d\zeta\]
\[\sup_{|\zeta-a|=r(a)} \left|\frac{k!}{2\pi i}\oint_{|\zeta-a|=r(a)}\frac{f_n(\zeta)-f(\zeta)}{(\zeta - z)^{k+1}} d\zeta\right| \leq \frac{k!}{2\pi i} \cdot \frac{\epsilon}{(r-R)^{k+1}} \text{ т.е. } \rightrightarrows\]
\[\forall K \subset D, K \text{ -- компакт. Рассмотрим } \{U_{R(a)}(a)\}_{a \in K} \text{ -- открытое покрытие, выберем конечное}\]
\[\text{подпокрытие } K \subset \bigcup_{n=1}^N \overline{U_{R_n(a_n)}(a_n)} \text{, т.к. их конечно, то } f_n^{(k)} \overset{\overline{U_{R_n(a_n)}(a_n)}}{\rightrightarrows} f^{(k)} \blacktriangleleft\]

\textbfit{Теорема} (Рунге). Пусть $D -$ односвязная область. $f \in \mathcal{O}(D)$. Тогда $\exists P_n -$ последовательность многочленов: $\forall F \subset D, K -$ компакт $P_n \overset{K}{\rightrightarrows} f$.

\textbfit{Пр.} $f(z) = \frac{1}{z}, \ \mathbb{C} \setminus \{0\}$
\[\oint_{|z|=1} P_n(z) dz = 0 \to \oint_{|z|=1} \frac{1}{z} dz = 2\pi i \Rightarrow \nexists P_n\]
